% Self-contained mathematical notes for incorporation into the main manuscript.
% No preamble; requires amsmath, amssymb, amsthm.  Prepared independently.

\section{A positive spectral representation of the entire remainder}
\label{sec:spectral-notes}

Put
\[
 a_j=\int_j^{j+1}\log x\,dx\quad(j\ge1),\qquad
 g_0=1,\qquad g_n=\sum_{j=1}^n a_jg_{n-j},
\]
and write $A(z)=\sum_{j\ge1}a_jz^j$ and
$G(z)=\sum_{n\ge0}g_nz^n=1/(1-A(z))$.  Let $\rho\in(0,1)$
be the unique solution of $A(\rho)=1$, put $w_0=\rho^{-1}$, and set
\[
 \gamma=\frac1{\rho A'(\rho)},\qquad
 h(t)=\frac{(1-t)^2}{\log^2t}\quad(0<t<1),\qquad c=2\log2.
\]
The continuous extensions are $h(0)=0$ and $h(1)=1$.  The constant
$\gamma$ is the residue amplitude; Euler's constant is denoted below
by $\gamma_{\!E}$.

\subsection{Elementary integral identities}

The identity
\[
 \frac{1-t}{-\log t}=\int_0^1t^u\,du
\]
gives the useful representation
\begin{equation}
 h(t)=\int_0^1\int_0^1t^{u+v}\,du\,dv.
 \label{eq:spec-h-mixture}
\end{equation}
In particular, Jensen's inequality and $0<t<1$ imply
\begin{equation}
 t\le h(t)\le1.
 \label{eq:spec-h-bounds}
\end{equation}
Direct differentiation of the formula for $a_j$, followed by the
elementary identity $\log(v/u)=\int_0^\infty(e^{-ux}-e^{-vx})\,dx/x$,
gives
\begin{equation}
 a_{j+1}-a_j=\int_0^1t^{j-1}h(t)\,dt\qquad(j\ge1).
 \label{eq:spec-differences}
\end{equation}
For completeness, one may also integrate the right side of
\eqref{eq:spec-h-mixture} first in $t$: it equals
$\int_0^1\int_0^1(j+u+v)^{-1}\,du\,dv$, and the two elementary
integrations yield
$(j+2)\log(j+2)-2(j+1)\log(j+1)+j\log j$.
Consequently
\[
 (1-z)A(z)=a_1z+z^2\int_0^1\frac{h(t)}{1-tz}\,dt,
 \qquad a_1=c-1.
\]
Define, off $[0,1]$,
\begin{equation}
 J(w)=w-c-\int_0^1\frac{h(t)}{w-t}\,dt.
 \label{eq:spec-J}
\end{equation}
Then initially for large $|w|$, and thereafter by analytic continuation,
\begin{equation}
 G(1/w)=\frac{w-1}{J(w)}.
 \label{eq:spec-G-J}
\end{equation}
Finally,
\begin{equation}
 \int_0^1\frac{h(t)}t\,dt
 =\int_0^\infty\frac{(1-e^{-x})^2}{x^2}\,dx=2\log2=c.
 \label{eq:spec-zeroidentity}
\end{equation}
Integration by parts reduces the last integral to
$2\int_0^\infty(e^{-x}-e^{-2x})\,dx/x=2\log2$.

\subsection{A bounded self-adjoint operator}

Work in the Hilbert space
$\mathcal H=\mathbb C\oplus L^2((0,1),h(t)\,dt)$ and define
\begin{equation}
 B(x,f)=\left(cx+\int_0^1f(t)h(t)\,dt,\ x+tf(t)\right).
 \label{eq:spec-operator}
\end{equation}
This is a bounded self-adjoint operator: multiplication by $t$ has
norm at most one and the off-diagonal map is the bounded rank-one
map induced by the constant function.  Moreover,
\begin{equation}
 \langle B(x,f),(x,f)\rangle
 =\int_0^1t h(t)\left|f(t)+\frac{x}{t}\right|^2\,dt\ge0.
 \label{eq:spec-positive}
\end{equation}
The displayed square is integrable because
$\int h(t)/t\,dt=c<\infty$; expanding it gives the preceding
quadratic form exactly.  Thus $B$ is positive.

Let $e_0=(1,0)$.  Solving $(w-B)(x,f)=e_0$ by its second component
gives $f(t)=x/(w-t)$ and hence
\begin{equation}
 \langle(w-B)^{-1}e_0,e_0\rangle=\frac1{J(w)}.
 \label{eq:spec-resolvent}
\end{equation}
By the bounded self-adjoint spectral theorem \cite{Teschl} there is a positive probability measure
$\sigma$, supported in the bounded spectrum of $B$, such that
\begin{equation}
 \frac1{J(w)}=\int\frac{\sigma(dt)}{w-t}.
 \label{eq:spec-measure}
\end{equation}

There is exactly one spectral point above one.  Indeed, on $(1,\infty)$,
\[
 J'(w)=1+\int_0^1\frac{h(t)}{(w-t)^2}\,dt>0,
\]
whereas $J(w)\to-\infty$ as $w\downarrow1$ and $J(w)\to\infty$
as $w\to\infty$.  The unique zero is $w_0=\rho^{-1}$.  The
Schur-complement calculation shows that every real $w\notin[0,1]$
with $J(w)\ne0$ is in the resolvent set.  Positivity excludes negative
spectrum.  Thus
\[
 \operatorname{supp}\sigma\subseteq[0,1]\cup\{w_0\}.
\]
The mass at $w_0$ is
\begin{equation}
 b_0=\sigma(\{w_0\})=\frac1{J'(w_0)},\qquad
 \gamma=\frac{w_0-1}{w_0J'(w_0)}=(1-\rho)b_0.
 \label{eq:spec-atom}
\end{equation}
The final equality can either be read from the residue in
\eqref{eq:spec-G-J} or checked by differentiating that identity.

\subsection{No hidden singular measure}

For $0<t<1$, put
\begin{align}
 R(t)&=t-c-\operatorname{PV}\int_0^1\frac{h(s)}{t-s}\,ds,\\
 J_+(t)&=R(t)+i\pi h(t).
 \label{eq:spec-boundary}
\end{align}
The sign of the imaginary part follows from the upper-half-plane
boundary value of $(w-s)^{-1}$.  Since $h$ is smooth and positive
on compact subintervals of $(0,1)$, the principal value and the
boundary values are continuous there; the convergence from the upper
half-plane is uniform on every smaller compact interval.  Because
$|J_+(t)|\ge\pi h(t)>0$, Stieltjes inversion applied on these
intervals proves absolute continuity, with density
\begin{equation}
 \frac{d\sigma}{dt}=\frac{h(t)}{|J_+(t)|^2}
 \quad(0<t<1).
 \label{eq:spec-density}
\end{equation}
For clarity, this also excludes singular continuous mass in the open
interval: inversion gives the measure of each compact interval as the
integral of the continuous boundary density; such intervals exhaust
$(0,1)$.

Only endpoint atoms remain to be checked.  From
\eqref{eq:spec-zeroidentity}, for $\varepsilon>0$,
\[
 J(-\varepsilon)
 =-\varepsilon\left(1+\int_0^1
            \frac{h(t)}{t(t+\varepsilon)}\,dt\right).
\]
Since $\int_0^1h(t)/t^2\,dt=\infty$,
\begin{equation}
 \sigma(\{0\})
 =\lim_{\varepsilon\downarrow0}\frac{-\varepsilon}{J(-\varepsilon)}=0.
 \label{eq:spec-nozero}
\end{equation}
The divergence follows immediately from
$h(t)\sim1/\log^2t$ at zero.  At the other endpoint,
$J(1+\varepsilon)\to-\infty$, so
\begin{equation}
 \sigma(\{1\})
 =\lim_{\varepsilon\downarrow0}\frac{\varepsilon}{J(1+\varepsilon)}=0.
 \label{eq:spec-noone}
\end{equation}
The isolated atom at $w_0$ contributes zero to this latter limit.
These arguments account for the entire measure:
\begin{equation}
 \sigma=b_0\delta_{w_0}
 +\frac{h(t)}{|J_+(t)|^2}\,\mathbf1_{(0,1)}(t)\,dt.
 \label{eq:spec-fullmeasure}
\end{equation}

\subsection{Exact positive remainder and its consequences}

Since $\sigma$ has total mass one, equations
\eqref{eq:spec-G-J} and \eqref{eq:spec-measure} give
\[
 G(1/w)-1=(w-1)\int\frac{\sigma(dt)}{w-t}-1
         =\int\frac{t-1}{w-t}\,\sigma(dt).
\]
Expanding at infinity is justified by bounded spectral support.
For every integer $n\ge1$,
\begin{equation}
 \boxed{\quad
 E_n:=\gamma\rho^{-n}-g_n
     =\int_0^1t^{n-1}\ell(t)\,dt>0,
 \qquad
 \ell(t)=\frac{(1-t)h(t)}{|J_+(t)|^2}.
 \quad}
 \label{eq:spec-exact-remainder}
\end{equation}
This indexing is essential.  The expression
$\int t^nk(t)\,dt$ with $k=\ell/t$ is equivalent for $n\ge1$,
but it must not be extended to $n=0$: the corresponding density is
not integrable at zero.  Indeed $J(-\varepsilon)\to0^-$ by dominated
convergence in \eqref{eq:spec-J} and \eqref{eq:spec-zeroidentity}, so
$G(-1/\varepsilon)-1\to+\infty$ in
\eqref{eq:spec-G-J}.  Its Cauchy representation is a bounded isolated
pole term plus $\int_0^1\ell(t)/(t+\varepsilon)\,dt$; monotone
convergence proves $\int_0^1\ell(t)/t\,dt=\infty$.
The exact generating-function identity is
\begin{equation}
 G(z)=1+\frac{\gamma\rho^{-1}z}{1-z/\rho}
             -z\int_0^1\frac{\ell(t)}{1-tz}\,dt.
 \label{eq:spec-regularized-generating}
\end{equation}

The representation proves strict complete monotonicity:
\begin{equation}
 (-1)^k\Delta^k E_n
 =\int_0^1t^{n-1}(1-t)^k\ell(t)\,dt>0
 \qquad(n\ge1,\ k\ge0),
 \label{eq:spec-complete-monotone}
\end{equation}
where $\Delta E_n=E_{n+1}-E_n$.  Every finite Hankel matrix
$(E_{r+i+j})_{0\le i,j\le m}$ is positive definite when $r\ge1$:
for any nonzero real polynomial $P(t)=\sum_{i=0}^mc_it^i$,
\[
 \sum_{i,j=0}^mc_ic_jE_{r+i+j}
 =\int_0^1t^{r-1}P(t)^2\ell(t)\,dt>0.
\]
In particular $E_{n+1}^2<E_nE_{n+2}$.  These are exact finite-$n$
statements, independent of any asymptotic expansion.

Several exact sum rules are also available.  Comparing the coefficient
of $1/w$ gives
\begin{equation}
 \int_0^1\ell(t)\,dt=\gamma\rho^{-1}-a_1<1.
 \label{eq:spec-massell}
\end{equation}
Monotone convergence and the total mass of $\sigma$ give
\begin{equation}
 \sum_{n=1}^\infty E_n
 =\int_0^1\frac{\ell(t)}{1-t}\,dt
 =1-b_0=1-\frac\gamma{1-\rho}.
 \label{eq:spec-sum0}
\end{equation}
Letting $w\downarrow1$ in \eqref{eq:spec-measure} shows that
$\int_0^1(1-t)^{-1}\,\sigma(dt)=b_0/(w_0-1)$: the continuous
part increases monotonically and the left side $1/J(w)$ tends to zero,
while the isolated-pole term tends to $-b_0/(w_0-1)$.  Consequently
\begin{equation}
 \sum_{n=1}^\infty nE_n
 =\int_0^1\frac{\ell(t)}{(1-t)^2}\,dt
 =\frac{b_0}{w_0-1}
 =\frac{\gamma\rho}{(1-\rho)^2}.
 \label{eq:spec-sum1}
\end{equation}

\section{Explicit finite-index bounds and the sharp leading scale}
\label{sec:spectral-bounds-notes}

We next prove actual numerical constants without needing certified
quadrature for the spectral density.  They are deliberately conservative;
the exact integral \eqref{eq:spec-exact-remainder} supports much tighter
numerical enclosures.

\subsection{The logarithmic endpoint constant}

Define
\[
 C=\int_0^1\frac{h(s)-1}{1-s}\,ds.
\]
With $s=e^{-x}$ this becomes
\[
 C=\int_0^\infty\left(
       \frac{e^{-x}-e^{-2x}}{x^2}-\frac1{e^x-1}\right)dx.
\]
Writing $\operatorname{E}_1(a)=\int_a^\infty e^{-x}\,dx/x$,
integration by parts after a cutoff $\varepsilon$ gives
\[
 C=\lim_{\varepsilon\downarrow0}
 \left\{
 \frac{e^{-\varepsilon}-e^{-2\varepsilon}}\varepsilon
 -\operatorname{E}_1(\varepsilon)
 +2\operatorname{E}_1(2\varepsilon)
 +\log(1-e^{-\varepsilon})\right\}.
\]
The standard integral definition
$\gamma_{\!E}=-\lim_{\varepsilon\downarrow0}
 (\operatorname{E}_1(\varepsilon)+\log\varepsilon)$ therefore yields
\begin{equation}
 C=1-2\log2-\gamma_{\!E}.
 \label{eq:spec-Cconstant}
\end{equation}
Only $0<\gamma_{\!E}<1$ is needed for the explicit bounds below.
These inequalities follow, for example, from
\[
 \gamma_{\!E}=\int_0^1\frac{1-e^{-x}}x\,dx
                  -\int_1^\infty\frac{e^{-x}}x\,dx.
\]

Put $r(s)=(h(s)-1)/(1-s)$ and extend it at one by $r(1)=-1$.
By \eqref{eq:spec-h-bounds}, $-1\le r(s)\le0$.  If $t=1-u$,
then a direct decomposition $h(s)=1+(1-s)r(s)$ gives
\begin{equation}
 R(1-u)=\log u+\gamma_{\!E}+\eta(u),\qquad
 \eta(u)=-u-\log(1-u)
       -u\operatorname{PV}\int_0^1\frac{r(s)}{1-u-s}\,ds.
 \label{eq:spec-etaidentity}
\end{equation}
Here $R$ is the real part in \eqref{eq:spec-boundary}.

We record explicit estimates for the last principal value.  From
\eqref{eq:spec-h-mixture}, for $s\in[1/4,1]$,
$|h''(s)|\le32$: indeed $|q(q-1)|\le2$ for $0\le q\le2$ and
$s^{q-2}\le16$.  Since $r$ is the negative average of $h'$ over
$[s,1]$, this implies $|r'(s)|\le16$ on that interval.  For
$t\in[1/2,1)$,
\[
 \operatorname{PV}\int_0^1\frac{r(s)}{t-s}\,ds
 =r(t)\log\frac{t}{1-t}
       +\int_0^1\frac{r(s)-r(t)}{t-s}\,ds.
\]
The integral over $[0,1/4]$ has absolute value at most one; the
remaining part has absolute value at most $16(3/4)=12$.  Therefore,
with $L=\log(1/u)$ and $0<u\le1/2$,
\begin{equation}
 |\eta(u)|\le u(L+14).
 \label{eq:spec-eta-bound}
\end{equation}
We used $0\le-u-\log(1-u)\le u^2$ for $u\le1/2$.
In particular,
\begin{equation}
 J_+(1-u)=\log u+\gamma_{\!E}+i\pi
                  +O(u\log(1/u)),
 \qquad
 \ell(1-u)\sim\frac{u}{\log^2(1/u)}.
 \label{eq:spec-local-equivalence}
\end{equation}

\subsection{Two-sided density bounds with actual constants}

For $0<u\le e^{-4}$, $L\ge4$ and
$u(L+14)\le18e^{-4}<1/3$.  Using $0<\gamma_{\!E}<1$,
\eqref{eq:spec-eta-bound} implies
\[
 \frac23L\le |R(1-u)|\le\frac43L.
\]
Also $h(1-u)\ge1-u>1/2$, $h(1-u)\le1$, and
$\pi h(1-u)\le\pi<L$.  Hence
\begin{equation}
 \boxed{\qquad
 \frac9{50}\frac{u}{\log^2(1/u)}
 \le\ell(1-u)\le
 \frac94\frac{u}{\log^2(1/u)}
 \quad(0<u\le e^{-4}).
 \qquad}
 \label{eq:spec-explicit-density}
\end{equation}

\subsection{A fully explicit finite-index theorem}

For every integer $n\ge4096$,
\begin{equation}
 \boxed{\qquad
 \frac{1}{200\,n(n+1)\log^2n}
 \le \gamma\rho^{-n}-g_n
 \le\frac{10}{n(n+1)\log^2n}.
 \qquad}
 \label{eq:spec-explicit-n}
\end{equation}
To prove the upper bound, split the integral in
\eqref{eq:spec-exact-remainder} at $u=n^{-1/2}\le e^{-4}$.
On the part $u\le n^{-1/2}$, $\log(1/u)\ge\tfrac12\log n$, so
\[
 \int_0^{n^{-1/2}}(1-u)^{n-1}\ell(1-u)\,du
 \le\frac9{\log^2n}\int_0^1(1-u)^{n-1}u\,du
 =\frac9{n(n+1)\log^2n}.
\]
On the complement use $\int\ell\le1$, a consequence of the
probability measure $\sigma$:
\[
 \int_{n^{-1/2}}^1(1-u)^{n-1}\ell(1-u)\,du
 \le e^{-\sqrt n+1/\sqrt n}
 \le\frac1{n(n+1)\log^2n}.
\]
For the last elementary inequality set $r=\sqrt n\ge64$.  The
difference between the logarithms of the two sides, with signs reversed,
is
\[
 r-r^{-1}-4\log r-\log(1+r^{-2})-2\log(2\log r).
\]
It is positive at $r=64$ (use $\log64<5$ and $\log10<3$), and its
derivative is positive for $r\ge64$.

For the lower bound restrict to $1/(2n)\le u\le1/n$.  On that
interval, $(1-u)^{n-1}\ge e^{-1}>1/3$ and
$\log(1/u)\le\log(2n)\le2\log n$.  Equation
\eqref{eq:spec-explicit-density} gives
\[
 E_n\ge\frac{9}{600\log^2n}
               \int_{1/(2n)}^{1/n}u\,du
      =\frac9{1600\,n^2\log^2n}
      \ge\frac1{200\,n(n+1)\log^2n}.
\]
No numerical root approximation, numerical principal value, or numerical
integration is used in \eqref{eq:spec-explicit-n}.

A weaker bound valid at every $n\ge2$ follows from
$|J_+(t)|^2\ge\pi^2h(t)^2$ and $h(t)\ge t$:
\begin{equation}
 0<E_n\le\frac1{\pi^2 n(n-1)}\qquad(n\ge2).
 \label{eq:spec-global-simple}
\end{equation}

\subsection{Sharp asymptotic constant}

The same estimates give
\begin{equation}
 \boxed{\qquad
 E_n\sim\frac1{n^2\log^2n}.
 \qquad}
 \label{eq:spec-leading}
\end{equation}
Here are details avoiding an implicit Tauberian hypothesis.  The part
$u\ge n^{-1/2}$ is negligible after multiplication by $n^2\log^2n$,
as just shown.  In the remaining integral put $v=nu$.  For each fixed
$v>0$, \eqref{eq:spec-local-equivalence} gives
\[
 n^2\log^2n\,(1-v/n)^{n-1}\ell(1-v/n)\,\frac{dv}{n}
 \longrightarrow ve^{-v}\,dv.
\]
For $v\le\sqrt n$, \eqref{eq:spec-explicit-density} bounds the
integrand by $9ve^{-v/2}$ for all sufficiently large $n$.  Dominated
convergence applies, and $\int_0^\infty ve^{-v}\,dv=1$.

Thus the error has a fixed negative sign in $g_n-\gamma\rho^{-n}$,
its absolute order $n^{-2}(\log n)^{-2}$ cannot be improved, and its
entire finite-index shape is constrained by a positive Hausdorff moment
representation.
